% =============================================================================
% Chapter 3 — The Vector Table ($0200–$027F)
% =============================================================================
\chapter{The Vector Table}

\epigraph{``Never hardcode what you can look up.  Addresses change; indirection
is forever.''}{\textit{--- Pragmatic advice from a systems programmer, ca.\ 1985}}

\section*{Introduction}

When the NovaVM boots, the monitor code at \$FF80 writes a table of 16-bit
addresses into the bottom of page 2.  Each entry holds the base address of a
major hardware device --- the graphics controller, the SID chip, the file
system, the network interface, and so on.  Programs that read these vectors
instead of hardcoding device addresses will continue to work even if the
hardware map changes in a future NovaVM version.

The vector table is ordinary RAM.  The system writes it once at boot, and
programs are free to read it at any time.  Writing to the table is technically
possible but strongly discouraged --- altering a vector will not move the
hardware, it will just confuse any code that reads the table afterward.

The 128-byte region from \$0200 to \$027F is reserved for the system vector
table.  The first 24 bytes (\$0200--\$0217) contain twelve 16-bit vectors,
stored in standard 6502 little-endian format (low byte at the even address,
high byte at the odd address).  The remaining 104 bytes (\$0218--\$027F) are
currently reserved.
\label{addr:VecTable}%
\index{\$0200@{\texttt{\$0200--\$027F}} --- System Vector Table}%

% ═══════════════════════════════════════════════════════════════════════════════
\section{Vector Entries}
% ═══════════════════════════════════════════════════════════════════════════════

Each vector is a two-byte pointer to the base register of a hardware
controller.  Use \texttt{DEEK} to read a vector as a 16-bit address:

\begin{lstlisting}
PRINT "$";HEX$(DEEK($0200))
\end{lstlisting}

This prints \texttt{\$A000} --- the VGC base address.

% ── VgcBase ──────────────────────────────────────────────────────────────────
\mementry{0200}{VgcBase (low)}{R/W}{\$00}

\mementry{0201}{VgcBase (high)}{R/W}{\$A0}

Points to \$A000, the base of the Virtual Graphics Controller register block.
The VGC core registers (mode, colours, cursor, scroll, collision) begin here.
See Chapter~\ref{addr:A000} for full VGC documentation.

\begin{lstlisting}
VB = DEEK($0200)
POKE VB+1, 6 : REM set background to blue
\end{lstlisting}

% ── RegCmdAddr ───────────────────────────────────────────────────────────────
\mementry{0202}{RegCmdAddr (low)}{R/W}{\$10}

\mementry{0203}{RegCmdAddr (high)}{R/W}{\$A0}

Points to \$A010, the VGC command register.  Writing a command byte here
triggers immediate execution of graphics primitives (plot, line, circle, fill)
and sprite operations.  Parameters are loaded into \$A011--\$A01E before the
command write.  See Chapter~\ref{addr:A000} for the complete command list.

% ── VramPortBase ─────────────────────────────────────────────────────────────
\mementry{0204}{VramPortBase (low)}{R/W}{\$E0}

\mementry{0205}{VramPortBase (high)}{R/W}{\$A0}

Points to \$A0E0, the VDC-style VRAM port base. Character RAM, Color RAM,
graphics bitmap memory, sprite shape RAM, and tile data RAM are selected by
plane number through this port.
See Chapter~\ref{sec:charram} for screen memory details.

\begin{lstlisting}
VP = DEEK($0204)
POKE VP, 1 : REM select character plane
\end{lstlisting}

% ── VramData ─────────────────────────────────────────────────────────────────
\mementry{0206}{VramData (low)}{R/W}{\$E3}

\mementry{0207}{VramData (high)}{R/W}{\$A0}

Points to \$A0E3, the VRAM data register. This shortcut is useful for tight
screen-fill loops after the caller has selected a plane and starting address
through the VRAM port.

% ── SidBase ──────────────────────────────────────────────────────────────────
\mementry{0208}{SidBase (low)}{R/W}{\$00}

\mementry{0209}{SidBase (high)}{R/W}{\$D4}

Points to \$D400, the base of SID chip 1.  The SID provides three analogue
synthesis voices with ADSR envelopes, ring modulation, and programmable
filters.  See Chapter~\ref{addr:D400} for register details.

% ── FileIoBase ───────────────────────────────────────────────────────────────
\mementry{020A}{FileIoBase (low)}{R/W}{\$A0}

\mementry{020B}{FileIoBase (high)}{R/W}{\$B9}

Points to \$B9A0, the File I/O controller.  This device handles loading and
saving BASIC programs, directory listing, graphics save/load, SID file
playback, and music sequencing.  See Chapter~\ref{addr:B9A0}.

% ── XmcBase ──────────────────────────────────────────────────────────────────
\mementry{020C}{XmcBase (low)}{R/W}{\$00}

\mementry{020D}{XmcBase (high)}{R/W}{\$BA}

Points to \$BA00, the Expansion Memory Controller.  The XMC provides access
to 512\,KB of expansion RAM through a register-based command interface and
four 256-byte memory windows at \$BC00--\$BFFF.
See Chapter~\ref{addr:BA00}.

% ── TimerBase ────────────────────────────────────────────────────────────────
\mementry{020E}{TimerBase (low)}{R/W}{\$40}

\mementry{020F}{TimerBase (high)}{R/W}{\$BA}

Points to \$BA40, the Timer controller.  A simple programmable interval timer
that generates IRQ interrupts at a configurable rate, useful for game loops
and real-time tasks.

% ── Sid2Base ─────────────────────────────────────────────────────────────────
\mementry{0210}{Sid2Base (low)}{R/W}{\$20}

\mementry{0211}{Sid2Base (high)}{R/W}{\$D4}

Points to \$D420, the base of SID chip 2.  A second complete SID with three
additional synthesis voices, giving six SID voices total.
See Chapter~\ref{addr:D400}.

% ── NicBase ──────────────────────────────────────────────────────────────────
\mementry{0212}{NicBase (low)}{R/W}{\$00}

\mementry{0213}{NicBase (high)}{R/W}{\$A1}

Points to \$A100, the Network Interface Controller.  The NIC provides four
TCP connection slots for client and server networking with DMA message
transfers and optional per-slot IRQ notification.
See Chapter~\ref{addr:A100}.

% ── DmaBase ──────────────────────────────────────────────────────────────────
\mementry{0214}{DmaBase (low)}{R/W}{\$63}

\mementry{0215}{DmaBase (high)}{R/W}{\$BA}

Points to \$BA63, the DMA controller.  High-speed block transfers between
CPU RAM, VGC memory spaces (character, colour, graphics, sprite), and
expansion RAM --- all without tying up the CPU.

% ── BlitterBase ──────────────────────────────────────────────────────────────
\mementry{0216}{BlitterBase (low)}{R/W}{\$83}

\mementry{0217}{BlitterBase (high)}{R/W}{\$BA}

Points to \$BA83, the Blitter controller.  A 2D block-transfer engine that
supports rectangular copy, fill, and colour-key transparency operations
across all memory spaces.

% ═══════════════════════════════════════════════════════════════════════════════
\section{Reserved Range (\$0218--\$027F)}
% ═══════════════════════════════════════════════════════════════════════════════

The 104 bytes from \$0218 to \$027F are currently unused.  They are ordinary
RAM and could be used by your own programs for variables or small data tables.
However, future NovaVM versions may add new hardware vectors here, so it is
best to treat this range as reserved if you want your programs to remain
forward-compatible.

% ═══════════════════════════════════════════════════════════════════════════════
\section{Reading the Complete Vector Table}
% ═══════════════════════════════════════════════════════════════════════════════

The following program dumps all twelve vectors:

\begin{lstlisting}
10 REM dump the vector table
20 FOR I = $0200 TO $0216 STEP 2
30   PRINT "$"; HEX$(I); " -> $"; HEX$(DEEK(I))
40 NEXT I
\end{lstlisting}

On a standard NovaVM, this produces:

\begin{lstlisting}
$0200 -> $A000
$0202 -> $A010
$0204 -> $A0E0
$0206 -> $A0E3
$0208 -> $D400
$020A -> $B9A0
$020C -> $BA00
$020E -> $BA40
$0210 -> $D420
$0212 -> $A100
$0214 -> $BA63
$0216 -> $BA83
\end{lstlisting}

\begin{notebox}
Programs should read the vector table rather than hardcoding device addresses.
This ensures forward compatibility if hardware base addresses change in future
NovaVM versions.  For example, instead of writing directly to \$A001, use:
\begin{lstlisting}
VB = DEEK($0200)
POKE VB+1, 6
\end{lstlisting}
\end{notebox}

% ═══════════════════════════════════════════════════════════════════════════════
\section*{Summary}
% ═══════════════════════════════════════════════════════════════════════════════

\begin{center}
\small
\begin{tabular}{@{}llll@{}}
\toprule
\textbf{Address} & \textbf{Vector Name} & \textbf{Points To} & \textbf{Device} \\
\midrule
\$0200 & VgcBase      & \$A000 & VGC core registers \\
\$0202 & RegCmdAddr   & \$A010 & VGC command register \\
\$0204 & VramPortBase & \$A0E0 & VRAM port base \\
\$0206 & VramData     & \$A0E3 & VRAM data register \\
\$0208 & SidBase      & \$D400 & SID chip 1 \\
\$020A & FileIoBase   & \$B9A0 & File I/O controller \\
\$020C & XmcBase      & \$BA00 & Expansion memory controller \\
\$020E & TimerBase    & \$BA40 & Timer controller \\
\$0210 & Sid2Base     & \$D420 & SID chip 2 \\
\$0212 & NicBase      & \$A100 & Network interface controller \\
\$0214 & DmaBase      & \$BA63 & DMA controller \\
\$0216 & BlitterBase  & \$BA83 & Blitter controller \\
\midrule
\$0218--\$027F & \multicolumn{3}{l}{\textit{Reserved for future use}} \\
\bottomrule
\end{tabular}
\end{center}
